% ---- pg-481 （书页 469） ----

$\pi_2$. 这样, $n_{2M}^{(H)}=1$.

期望真实失误率的估计值为
\[
\hat{E}(\mathrm{AER})=\frac{n_{1M}^{(H)}+n_{2M}^{(H)}}{n_1+n_2}=\frac{2+1}{3+3}=0.5
\]
这样, 我们看到表观失误率 $\mathrm{APER}=0.33$ 是一个乐观的效果量度. 当然, 实践中的样本容量要比我们这里所考虑的容量大, 因而 $\mathrm{APER}$ 与 $\hat{E}(\mathrm{AER})$ 之间的差别没有本例中那么大. $\blacksquare$

有兴趣进一步了解分类失误率的估计方法的读者, 可参考文献[23].

下一个例子将说明当两个总体的判别方差不同时可能产生的一个困难.

\noindent\textbf{例 11.8}\quad (阿拉斯加产鲑鱼和加拿大产鲑鱼的分类)

捕捞鲑鱼对美国和加拿大两国来说都是一种有价值的资源. 由于这种资源有限, 因此必须有效地管理. 此外, 由于这个问题与两个国家有关, 必须公平地解决. 也就是说, 阿拉斯加的捕鱼者不能捕捞过多的加拿大产鲑鱼, 反之亦然.

这种鱼有一个奇怪的生命周期, 它们出生于淡水激流, 经过一到两年后游入大洋. 在海水中生活两年后, 它们又洄游到出生地, 产卵后死去. 在它们即将成熟, 在由大洋回游到淡水之前, 人们捕捞它们. 为了帮助规范捕捞, 在捕鱼季节, 管理人员从捕获的鱼中进行抽样检查, 以识别鱼是来自阿拉斯加还是来自加拿大产鲑鱼. 鱼鳞上的年轮可揭示产地. 典型的情况是: 阿拉斯加产鲑鱼的淡水生长年轮比加拿大产鲑鱼的要小. \hyperref[tab:11.2]{表 11.2} 给出了年轮的直径数据, 其中

\begin{align*}
X_1&=\text{第一年淡水生长的年轮直径(百分之一英寸)}\\
X_2&=\text{第一年海洋生长的年轮直径(百分之一英寸)}
\end{align*}

此外, 雌鱼取代码 1, 雄鱼则取代码 2.

\begin{mvtable}
\phantomsection\label{tab:11.2}\mvtabcap{表 11.2\quad 鲑鱼数据(年轮直径)}
\begin{tabular}{cccccc}
\toprule
\multicolumn{3}{c}{阿拉斯加} & \multicolumn{3}{c}{加拿大}\\
\cmidrule(lr){1-3}\cmidrule(lr){4-6}
性别 & 淡水 & 海洋 & 性别 & 淡水 & 海洋\\
\midrule
2 & 108 & 368 & 1 & 129 & 420\\
1 & 131 & 355 & 1 & 148 & 371\\
1 & 105 & 469 & 1 & 179 & 407\\
2 & 86 & 506 & 2 & 152 & 381\\
1 & 99 & 402 & 2 & 166 & 377\\
2 & 87 & 423 & 2 & 124 & 389\\
1 & 94 & 440 & 1 & 156 & 419\\
2 & 117 & 489 & 2 & 131 & 345\\
2 & 79 & 432 & 1 & 140 & 362\\
1 & 99 & 403 & 2 & 144 & 345\\
1 & 114 & 428 & 2 & 149 & 393\\
2 & 123 & 372 & 1 & 108 & 330\\
1 & 123 & 372 & 1 & 135 & 355\\
2 & 109 & 420 & 2 & 170 & 386\\
2 & 112 & 394 & 1 & 152 & 301\\
1 & 104 & 407 & 1 & 153 & 397\\
2 & 111 & 422 & 1 & 152 & 301\\
\bottomrule
\end{tabular}
\end{mvtable}
